\section{A nonconserved collision-compatible acquisition}
\label{sec:v16-physical}
The following example lies on the full equilibrium phase space of an
interacting hard-sphere system. It is neither a collision-free restriction
nor an observable constant along the flow. Its explicitly solvable
collective mode permits simultaneous evaluation of a nonzero generator
residual and of a marked decision deficiency. The orbit factors through total momentum, so collisions cancel from
this particular observable. The residual measures projection error,
not interaction strength. The collision-sensitive comparison is proved
in Section~\ref{sec:v17-collision}; the solvable collective calculation
is retained here as a separate exact benchmark.

Fix $N\ge2$ equal-mass hard spheres on the square flat torus of side
$\ell>0$, with positive admissible configuration volume. Velocities
have independent standard Gaussian coordinates under equilibrium $\mu$.
Let $U(t)$ be the full hard-sphere Koopman group, and allow the two
physical interventions $v\mapsto av$, $a\in\{-1,1\}$, before each
flight of length $\Delta$. Define the circle-valued coordinate and the
normalized total momentum
\[
 \vartheta=\frac{2\pi}{\ell}\sum_{i=1}^N q_{i,1}\pmod {2\pi},\qquad
 X=N^{-1/2}\sum_{i=1}^N v_{i,1},\qquad
 \omega=\frac{2\pi\sqrt N}{\ell},\qquad f=\cos\vartheta.
\]
The exponential of the displayed coordinate is well defined on the torus;
no choice of lifted particle positions enters $f$.

\begin{theorem}[Collective domain and controlled-word reduction]
\label{thm:v16-collective}
The coordinate $\vartheta$ is Haar uniform and independent of the
standard normal $X$. On the full-measure domain of the hard-sphere flow,
\begin{equation}\label{eq:v16-collective-orbit}
 U(t)f=\cos(\vartheta+\omega Xt),\qquad
 L^j f=\omega^jX^j\cos(\vartheta+j\pi/2),\qquad
 \|L^jf\|_2^2=\frac{\omega^{2j}(2j-1)!!}{2}.
\end{equation}
Here $f\in\bigcap_{j\ge1}D(L^j)$ and $(-1)!!=1$. In the normalized
one-dimensional trial space spanned by $\varphi=\sqrt2f$, the matrices
of the preceding residual theorem are
\[
 G=1,\qquad J=0,\qquad R=\omega^2>0,\qquad
 B_{-1}=B_1=1,\qquad S_{-1}=S_1=0.
\]
For a word $w=(a_1,\ldots,a_k)$ put
\[
 s(w)=\Delta\sum_{j=1}^k\prod_{i=1}^j a_i.
\]
Then $T_{a_1}\cdots T_{a_k}f=\cos(\vartheta+\omega Xs(w))$.
Exactly $k+1$ distinct observable functions occur among the $2^k$
words, indexed by $s/\Delta\in\{-k,-k+2,\ldots,k\}$.
\end{theorem}
\begin{proof}
Simultaneous translation of every particle by $r$ in the first torus
coordinate preserves the admissible configuration measure and shifts
$\vartheta$ by $2\pi Nr/\ell$. These shifts exhaust the circle, so its
push-forward measure is Haar. Configuration and velocity are independent,
and Gaussian addition proves the assertion about $X$.
Elastic collisions conserve total momentum and do not jump particle
positions. Between collisions the derivative of $\vartheta$ is
$\omega X$. This proves the orbit formula, including the periodic
identifications, on the classical almost-everywhere flow domain.
Difference quotients of that formula, dominated successively by
polynomials in $|X|$, converge in $L^2(\mu)$ since all Gaussian moments
are finite. Induction gives membership in every generator power domain
and the derivative formula. Independence and the uniform circle integral
of squared sine or cosine give the norm formula. This also establishes
collision compatibility in the actual generator domain, rather than
by applying a formal derivative to arbitrary boundary data.

The inner product of $\cos\vartheta$ and $X\sin\vartheta$ vanishes.
Thus $J=0$ and the squared residual norm of $L\varphi$ is $\omega^2$.
A velocity reversal fixes $\varphi$, giving $B_a=1$ and $S_a=0$.
After the $j$th intervention the total momentum is
$(a_1\cdots a_j)X$, so summing flight displacements proves the word
formula in the order $T_a=C_aU(\Delta)$. Every sequence of $k$ signs
can be the successive cumulative products, which gives all $k+1$
clock values. They are distinct as observable functions because
\[
 \|\cos(\vartheta+\omega Xs)-\cos(\vartheta+\omega Xt)\|_2^2
       =1-\exp\{-\omega^2(s-t)^2/2\}>0\quad(s\ne t).
\]
\end{proof}

\begin{theorem}[Explicit converging finite physical presentations]
\label{thm:v16-taylor}
Let $m\ge1$ and replace the mean at clock $s$ by
\[
 f_{m,s}=\sum_{j=0}^{m-1}\frac{s^j}{j!}L^jf.
\]
For $|s|\le k\Delta$ the mean error is bounded by
\begin{equation}\label{eq:v16-taylor-error}
 E_{m,k}=\frac{(k\Delta\omega)^m}{m!}
                  \sqrt{\frac{(2m-1)!!}{2}}.
\end{equation}
For Gaussian report noises $\sigma_k>0$ and preparations satisfying
$\sup_\theta\|p_\theta\|_2\le D$, the exact and polynomial
acquisitions have two-sided stateless marked feedback comparison error
at most
\begin{equation}\label{eq:v16-wordbound}
 \alpha_{m,T}=1\wedge\frac{D}{\sqrt{2\pi}}
             \sum_{k=1}^T\frac{\sqrt{k+1}}{\sigma_k}E_{m,k}.
\end{equation}
The same actual initial mark and every exposed independent selector are
preserved. The estimate is uniform over all consumer register widths,
and $\alpha_{m,T}\to0$ at every fixed finite horizon.
The physical trial functions and all their generator and intervention
Gram entries are explicitly integrable against equilibrium.
\end{theorem}
\begin{proof}
The integral Taylor remainder for the unitary orbit on $D(L^m)$ has
norm at most $|s|^m\|L^mf\|_2/m!$, also for negative $s$.
Theorem~\ref{thm:v16-collective} gives \eqref{eq:v16-taylor-error}.
Its successive ratio is
\[
 \frac{E_{m+1,k}}{E_{m,k}}
       =k\Delta\omega\frac{\sqrt{2m+1}}{m+1}\longrightarrow0,
\]
which is an effective bound for the convergence claimed here.

Couple the two acquisitions from the same initial $x,W$ and the same
independent selector. Until the first unequal report, their controls
agree. Two unit-dimensional Gaussians with the same variance have
variation distance at most their mean difference divided by
$\sigma_k\sqrt{2\pi}$; this follows by integrating the derivative of
the translated Gaussian density. At fixed $x$, the probabilities of
all matched prefixes sum to at most one. Keep those probabilities
until bounding the mismatch by the largest mean error at that $x$.
There are only $k+1$ different errors at time $k$, even though there
are $2^k$ words. Cauchy--Schwarz gives
\[
 \int p_\theta(x)\max_s|e_{m,s}(x)|\,d\mu(x)
 \le D\Bigl(\sum_s\|e_{m,s}\|_2^2\Bigr)^{1/2}
 \le D\sqrt{k+1}E_{m,k}.
\]
Sum first-mismatch probabilities over $k$. The construction in both
directions uses the identity report interface, has no added persistent
state, and couples the actual mark, proving \eqref{eq:v16-wordbound}.
No bound on feedback-conditioned preparation densities was used.

The trial family consists of $X^j\cos\vartheta$ and
$X^j\sin\vartheta$ for $j<m$. Its inner products are products of
circle integrals and Gaussian moments. Generator entries multiply by
$\omega X$ and rotate sine and cosine; the interventions replace $X$
by $-X$. Every required finite entry is thus a polynomial in $\omega$
with rational coefficients before normalization. This supplies exact
symbolic expressions and arbitrarily tight rational interval enclosures.
The polynomial means may be unbounded in $X$; that is harmless for
\eqref{eq:v16-wordbound}. For finite-report integration one first
truncates $|X|$ with its explicit Gaussian tail, then integrates bounded
Gaussian report probabilities over the resulting compact cylinder.
For the equilibrium and mark used below, the only mark boundaries are
$\vartheta=\pi/2,3\pi/2$, and can be split explicitly. On each piece
trigonometric, polynomial and Gaussian distribution functions have
computable uniform moduli. Rational upper and lower Riemann enclosures
therefore converge. This verifies the effective marked integration
hypothesis of the finite-presentation theorem for this class, without
an exact hard-sphere flow integration oracle.
\end{proof}

\begin{remark}
The signed clock is an index of physical mean functions, not an
unpriced extra register offered to a consumer. The width-uniform proof
uses equality of mean functions to reduce a bound; it does not allow
a finite-state simulator to remember its entire control word. Neither
the polynomial coefficient vector nor its real dimension is a finite
consumer-label count. The underlying microscopic system still has
collisions; solvability here comes from the conserved total momentum
that transports a \emph{nonconserved} position observable.
\end{remark}

\subsection{One actual mark, one nonzero residual, and a positive certificate}
Let $W=\operatorname{sgn}\cos\vartheta$ in the initial equilibrium
state; the zero set has measure zero. Both experiments have a singleton
initial report and allow $a\in\{-1,1\}$ before their single acquisition.
The source reports an independent standard Gaussian. The target reports
\[
 Y=\cos(\vartheta+a\tau X)+\xi,\qquad
 \tau=\omega\Delta,\qquad \xi\sim N(0,1)
\]
with noise independent of the actual initial mark and microscopic state.
Let $\Phi$ and $\phi$ denote the standard normal distribution and density.

\begin{theorem}[Exact marked physical deficiency]
\label{thm:v16-exactphysical}
Here the action precedes any informative report and its reversal
symmetry leaves the relevant marked law unchanged. In the visible
comparison the independent selector is retained jointly with the same
selector law on the target side. For every nonempty finite private,
hidden-selector, or visible-selector signature in this one-acquisition model, its intrinsic deficiency equals
\begin{equation}\label{eq:v16-delta}
 d(\tau)=\mathbb E\bigl[W\{\Phi(\cos(\vartheta+\tau X))-1/2\}\bigr]>0.
\end{equation}
It is attained by the stateless simulator that draws the target's
unconditional report marginal. If
\[
 c_k(\tau)=2^{-(2k+1)}\sum_{j=0}^{2k+1}\binom{2k+1}{j}
 \frac{2(-1)^{(|2k+1-2j|-1)/2}}{\pi|2k+1-2j|}
 e^{-(2k+1-2j)^2\tau^2/2},
\]
then
\begin{equation}\label{eq:v16-series}
 d(\tau)=\frac1{\sqrt{2\pi}}
      \sum_{k=0}^\infty\frac{(-1)^kc_k(\tau)}{2^kk!(2k+1)}.
\end{equation}
The error after $k=K$ is at most
$[\sqrt{2\pi}\,2^{K+1}(K+1)!(2K+3)]^{-1}$.
\end{theorem}
\begin{proof}
The action is chosen before any informative report. Reversal of the
independent symmetric $X$ does not change the joint law of $(W,Y)$,
including conditionally on any exposed independent seed. Set
$H_\tau(\theta)=\mathbb E\operatorname{sgn}\cos(\theta+\tau X)$.
This is the circle heat evolution of $\operatorname{sgn}\cos\theta$
at time $\tau^2/2$. It is strictly positive on $(-\pi/2,\pi/2)$ and
strictly negative on the opposite arc. To see the sign without assuming
it from a Fourier expansion, reflect the heat solution oddly across
both endpoints of the positive arc. Its restriction solves the heat
equation on that interval with zero boundary values and initial value
one. Equivalently it is the survival probability of Brownian motion
killed at those endpoints. This is positive in the interior at every
finite positive time. The odd reflected solution agrees with the
circle solution by uniqueness (the difference has zero initial data
and its squared integral has nonpositive derivative). At $\tau=0$
the same sign assertion is immediate.

The final angle is uniform. Reversibility of the Gaussian circle
convolution gives $\mathbb E[W\mid\vartheta+\tau X=\theta]
=H_\tau(\theta)$. Let $\nu$ be the unconditional density of $Y$.
The density of $(W,Y)$, minus $\tfrac12\nu(y)$ in its $w$th row, is
\[
 \frac{w}{4\pi}\int_0^{2\pi}H_\tau(\theta)
                                      \phi(y-\cos\theta)\,d\theta.
\]
Pair points $\theta$ and $\theta+\pi$. On the positive arc both
$H_\tau(\theta)$ and $\cos\theta$ are positive, and
$\phi(y-\cos\theta)-\phi(y+\cos\theta)$ has the sign of $y$.
Thus the density difference has the sign of $wy$, and is not zero.
Its positive event is $\{wy>0\}$. The excess probability of that event
is \eqref{eq:v16-delta}, proving positivity and giving a stateless upper
witness of that variation distance.

Every source simulator at the declared signature, conditionally on
its exposed seed, produces a report independent of the fair $W$.
For any such report law the event $\{wy>0\}$ has probability at most
$1/2$, including possible mass at zero. The target probability is
$1/2+d(\tau)$. This supplies the matching lower bound for private and
hidden randomization; the same bound in each visible fiber supplies
it for visible randomization. Mixtures of the two action choices do
not change this conclusion.

Finally, uniformly for $|z|\le1$,
\[
 \Phi(z)-1/2=\frac1{\sqrt{2\pi}}
       \sum_{k\ge0}\frac{(-1)^kz^{2k+1}}{2^kk!(2k+1)}.
\]
The absolute remainder is bounded by the next term, by the alternating
Taylor expansion of $e^{-u^2/2}$ followed by integration. Expand
$\cos^{2k+1}$ into its finite complex exponential sum. For odd $n$,
\[
 \mathbb E[W\cos(n\vartheta)]
       =\frac{2(-1)^{(|n|-1)/2}}{\pi|n|},\qquad
 \mathbb E e^{in\tau X}=e^{-n^2\tau^2/2}.
\]
The sine terms vanish. These identities give $c_k$ and the claimed
series and remainder, with no unproved exchange of a nonuniform sum.
\end{proof}

\begin{corollary}[A rational, nonvacuous collision certificate]
\label{cor:v16-numerical}
For $N=2$, $\ell=10$, any rational diameter in $(0,5)$ giving nonempty
interior, and $\Delta=1/2$, the same physical acquisition satisfies
\[
 \frac{205809122887}{10^{12}}<d(\tau)
            <\frac{205809122889}{10^{12}},\qquad
 R=\frac{2\pi^2}{25}>0.
\]
The degree-fifteen physical Taylor presentation has marked comparison
error less than $3\cdot10^{-11}$ for the single report.
\end{corollary}
\begin{proof}
Here $\tau^2/2=\pi^2/100$. Evaluate \eqref{eq:v16-series} through
$K=10$ with outward rational arithmetic and add its stated remainder.
For reproducibility the source supplies every arithmetic operation:
$\pi=16\arctan(1/5)-4\arctan(1/239)$ is bounded by alternating series;
square roots are enclosed by rational bisection; exponentials are
range-reduced to arguments in $[0,1]$, enclosed by alternating series,
and squared back with outward rounding. The resulting enclosure has
lower endpoint greater than the displayed lower rational and upper
endpoint smaller than the displayed upper rational. All comparisons
are integer comparisons. Applying \eqref{eq:v16-taylor-error} with
$m=16$, and replacing $\sqrt2$ in \eqref{eq:v16-wordbound} by the
larger rational $2$, yields the final strict bound. The generator
residual follows exactly from Theorem~\ref{thm:v16-collective}.
The supplied checker redoes these finite rational computations; the
analytic justification is the preceding domain, coupling and series
proofs, not the execution of that checker.
\end{proof}
